\section{ADR-007: File-backed JSON stores instead of a database}
\label{adr:007}

\subsection*{Status}

Accepted; in force at revision \texttt{\projectrevisionshort}.

\subsection*{Context}

Five kinds of state outlive a request: conversations, ingestion jobs, saved
insight reports, human answer feedback, and active-learning workflow state.

\subsection*{Decision}

Each is persisted as JSON files under a configurable directory beneath
\srcfile{.artifacts/}, one file per entity, through a small store class
(\texttt{JsonConversationStore}, \texttt{JsonJobStore},
\texttt{JsonInsightReportStore}, \texttt{JsonHumanFeedbackStore},
\texttt{JsonActiveLearningStateStore}). Records are re-validated against their
Pydantic model on load. No database, ORM, or migration tool is present in the
dependency set.

\subsection*{Consequences}

\begin{itemize}
  \item No service is required to run the API, and CI needs no database
    container.
  \item Stored state is inspectable and diffable with ordinary tools, which
    makes the workflow features easy to test.
  \item Each store is behind a class boundary, so replacing the backing store
    is a contained change.
  \item There is no transactionality and no cross-process locking. Concurrent
    writers to the same entity race on whole-file writes. The production
    Compose file runs multiple Gunicorn workers, which makes this reachable in
    that configuration.
  \item Listing is a directory scan, so it degrades linearly with the number of
    stored entities.
\end{itemize}

\subsection*{Alternatives}

SQLite, which would add transactionality at the cost of a schema and
migrations; or an external database, which would contradict \cref{adr:001}.
The repository records no evaluation of either.

\noevidence

\subsection*{Evidence}

\srcfile{src/feedback\_intelligence\_agent/memory.py}, \srcfile{jobs.py},
\srcfile{reports.py}, \srcfile{human\_feedback.py},
\srcfile{active\_learning.py};
\srcfile{src/feedback\_intelligence\_agent/config.py} (the five
\texttt{*\_store\_path} settings); \srcfile{pyproject.toml} (no database
dependency); \srcfile{deploy/docker-compose.prod.yml} (multiple workers).
